We separate \emph{reproduced results} (from prior work, replicated as
sanity checks) from \emph{planned evaluation} (the headline numbers
the v1.0 release targets). The v0.1 release ships the unit-test suite
and the four end-to-end examples; the benchmark harnesses described
below live under \texttt{benchmarks/} and are reproduced as part of
the v1.0 release process.

\subsection{Benchmarks}

\paragraph{AgentDojo (for \texttt{aisafepy.flow}).} AgentDojo
\citep{debenedetti2024agentdojo} measures attack success rate (ASR)
and utility on indirect-prompt-injection scenarios across a workspace
simulator. We target the canonical configuration: GPT-4o-class
baseline, the public workspace tools, the canonical
``important-message'' attack family. The success criterion is
ASR $\leq 5\%$ while preserving $\geq 90\%$ utility relative to the
unguarded baseline, matching the numbers reported by CaMeL on the
same benchmark \citep{camel2025}.

\paragraph{Streaming guard latency (for \texttt{aisafepy.stream}).}
The streaming-guard benchmark we ship measures macro-F1 at varying
token budgets and reports time-to-first-flag against a held-out
mixed corpus of jailbreak attempts and benign traffic. We target
macro-F1 $\geq$ that of the underlying classifier at $\leq$80 ms p95
added latency on a single A10 GPU. The cascade architecture is
designed so this is achievable without retraining: Tier 1 catches
deterministic patterns immediately, Tier 2 catches most policy
violations within $\sim$40 ms, Tier 3 fires only on uncertainty.

\paragraph{HarmBench (cross-module).} HarmBench
\citep{harmbench} measures direct content-policy violations across
domains. The cascade's Tier 2 plus Tier 3 should match the
underlying Llama Guard 4 score within $\pm 2$ F1 points while
reducing p95 latency by an order of magnitude through Tier 1
filtering and Tier 3 gating.

\paragraph{AILuminate (for \texttt{aisafepy.adapt}).} AILuminate
provides a sustained corpus of red-team scenarios. Our criterion is
not a static score but a velocity metric: time-to-fix from a
red-team finding to a deployed guard. The baseline is the manual
workflow (read JSONL, hand-write a Guardrails AI validator) which we
measure at roughly two engineer-days per cluster. The compiler's
target is $\leq$30 minutes from \texttt{compiler.compile()} to a
canary-deployed artifact.

\subsection{Microbenchmarks}

We run two microbenchmarks repeatedly in CI on Python 3.10/3.11/3.12:

\begin{enumerate}[leftmargin=2em, itemsep=0.25em]
  \item Tier-1 regex guard throughput on a synthetic 1MB streaming
    corpus. The current implementation, using the standard library
    \texttt{re} module, sustains $\geq$50,000 tokens per second per
    core. Switching to \texttt{re2} or \texttt{hyperscan} (optional
    dependencies) is benchmarked at $\geq$200,000 tokens per second
    per core on the same corpus.

  \item Adversarial false-positive resistance: we replay a corpus
    of short adversarial prompts known to drive content classifiers
    into high false-positive rates against a default
    \texttt{RegexGuard.compile\_pii()} configuration with the
    50,000-character input cap enabled. The cap prevents the
    DoS-on-guardrail attack from cleanly amplifying.
\end{enumerate}

\subsection{Reproducibility}

Each benchmark harness ships as a $\sim$200-LOC driver under
\texttt{benchmarks/} plus a YAML config that pins the model, the
dataset version, and the hardware. The full suite is intentionally
not part of the default \texttt{pytest} run because it requires GPU
access and downloaded model weights. The repository's GitHub Actions
workflow runs only the unit suite on CPU; a separate weekly workflow
runs the GPU benchmarks and publishes results to the
\texttt{benchmarks/results/} directory.

The numbers reported in the next revision of this paper will be
generated by these harnesses on commodity cloud GPUs (single A10
class) so that any reader with comparable hardware can reproduce
them. We will publish raw logs, intermediate artifacts, and
evaluation seeds for every reported value.

\subsection{Threat model and out-of-scope adversaries}

The evaluation explicitly excludes three threat models:

\begin{enumerate}[leftmargin=2em, itemsep=0.25em]
  \item \textbf{Sleeper-agent / deceptive alignment.} Anthropic's
    probe result \citep{anthropicprobes} is reproducible on toy
    backdoors and supports \texttt{stream.probes}, but no current
    technique reliably detects deception in frontier models. We do
    not claim coverage.
  \item \textbf{Side-channel attacks against model weights.}
    Activation-steering and probe-based defenses can themselves be
    evaded by an adversary who can fine-tune the underlying model.
    AIsafePy targets the deployment surface, not the weight-update
    surface.
  \item \textbf{Compromised execution environment.} The library
    assumes its own process is not compromised. A malicious host
    that intercepts the IFC interpreter's bytecode can bypass any
    runtime guarantee.
\end{enumerate}
